\section{Introduction}

GPU graph systems commonly organize work around a changing frontier and a
bulk-synchronous sequence of parallel edge operations.  Gunrock makes the
frontier explicit~\cite{wang2016gunrock}; Ligra separates sparse and dense
edge traversals~\cite{shun2013ligra}; and PowerGraph exposes a
gather--apply--scatter decomposition~\cite{gonzalez2012powergraph}.  GraphBLAS
instead specifies graph algorithms using sparse linear algebra over user
selected semirings~\cite{kepner2016graphblas}.  These systems demonstrate that
a small set of composable operations can cover many important algorithms, but
coverage by examples alone does not identify the exact expressive boundary.

\TA{} is intended to make that boundary a theorem.  Its central object is not
a materialized edge list or a particular kernel.  It is the denotation of an
ordered frontier traversal.  Expressions pull source, destination, and edge
attributes into an edge context; a map transforms each context; and a terminal
emits or reduces the resulting values.  Implementations may fuse these stages,
change materialization strategies, or select different parallel primitives as
long as they preserve the denotation.

The principal result is a characterization theorem, not a collection of
examples: for one arbitrary fixed scalar signature, every program in the
stated closed, typed, frontier-to-frontier \TA{} grammar has a
semantics-equivalent Monoidal Frontier BSP program, and conversely.  The proof
passes through an independently defined compositional syntax and a
sharing-aware single-push normal form.  It preserves complete finite runs and
empty-frontier termination, has factor-one syntax growth with exact
language-level work/depth/local-storage equations, covers dense and stable
sparse frontier selection, and separately equates the ordered observations of
emission and source/destination reduction.  These qualifications identify the
proved boundary without calling it unrestricted graph or Turing completeness.

This document has three deliberately separated levels of assurance:

\begin{enumerate}
  \item mathematical definitions and proofs presented here;
  \item kernel-checked counterparts of the language results and a typed
        abstract CubeCL resource semantics in Lean~4;
  \item differential value tests against the public
        \texttt{massively::graph} implementation, plus generated checks of the
        serialized abstract resource certificates.
\end{enumerate}

The third level is empirical and does not turn the Rust or CubeCL compiler into
part of the formal proof.  Conversely, the Lean development proves statements
about the mathematical language and abstract target without assuming that the
emitted CubeCL IR or implementation is correct.  Keeping these claims distinct
is essential to making later results auditable.
